\chapter{La continuità evolutiva dei bias: dal topo all'intelligenza collettiva}

\begin{flushright}
	\textit{``La cooperazione reciproca, non la lotta di ciascuno contro tutti, è stata l'arma principale nella battaglia per la sopravvivenza.''}\\[6pt]
	\textsc{\index[main]{Kropotkin, Pyotr}Pyotr Kropotkin}
\end{flushright}

\bigskip
\bigskip

C'è qualcosa di profondamente umiliante e al tempo stesso liberatorio nel guardare un \index[main]{topo}topo di laboratorio che sviluppa rituali superstiziosi, una colonia di \index[main]{formiche}formiche che risolve problemi complessi senza che nessuna formica capisca cosa sta facendo, o uno \index[main]{scimpanzé}scimpanzé che mostra la stessa \index[cognitive]{avversione alle perdite}avversione alle perdite di un trader di Wall Street. Non siamo soli nella nostra irrazionalità. È una storia di \index[main]{continuità evolutiva}continuità evolutiva che attraversa l'intero regno animale, dal più umile roditore fino a noi, presunti padroni del pianeta.

\section{Il topo filosofo non sopravvive}

Immaginate due topi. Il primo è quello che potremmo chiamare un \index[main]{topo filosofo}topo filosofo: quando sente un rumore si ferma, analizza, valuta le probabilità che sia un predatore oppure il vento. Il secondo è un \index[main]{topo paranoico}topo paranoico: scappa subito, sempre, senza pensarci due volte. Dopo un milione di anni di \index[main]{evoluzione}evoluzione, quale dei due ha più discendenti? La risposta è ovvia. Il topo paranoico spreca energie nel novantanove per cento dei casi, ma resta vivo per raccontarlo. Il topo filosofo ha ragione il novantanove per cento delle volte, ma quell'uno per cento di errore è fatale. Come dicono i biologi evoluzionisti, è meglio confondere cento volte un bastone per un serpente che confondere una sola volta un serpente per un bastone.\footnote{Haselton, M. G., e Nettle, D. (2006). The paranoid optimist: An integrative evolutionary model of cognitive biases. \textit{Personality and Social Psychology Review}, 10(1).}

Questo principio, che gli scienziati chiamano \index[main]{error management theory}teoria della gestione dell'errore, spiega perché i bias cognitivi non siano difetti ma caratteristiche. Nei laboratori i topi mostrano chiari segni di quello che negli umani chiamiamo \index[cognitive]{bias di conferma}bias di conferma: se hanno imparato che un certo suono precede il cibo, continueranno a cercarlo anche quando il suono non è più correlato al premio.\footnote{Rescorla, R. A. (1988). Pavlovian conditioning: It's not what you think it is. \textit{American Psychologist}, 43(3).} Ma il fenomeno diventa ancora più affascinante con i \index[main]{piccioni}piccioni. In uno degli esperimenti più celebri della psicologia, Burrhus Skinner scoprì che i piccioni sviluppano veri e propri rituali superstiziosi.\footnote{Skinner, B. F. (1948). Superstition in the pigeon. \textit{Journal of Experimental Psychology}, 38(2).} Se un piccione per caso stava girando su se stesso quando è arrivato del cibo distribuito a caso, continuerà a girare nella speranza che il rituale funzioni di nuovo. Uno faceva otto giri in senso antiorario tra un premio e l'altro, un altro sbatteva la testa contro un angolo della gabbia.

Vi suona familiare? Quanti di noi hanno una maglietta fortunata per gli esami? Quanti atleti seguono rituali elaborati prima di una gara? La \index[main]{superstizione}superstizione non è stupidità: è l'effetto collaterale di un sistema di apprendimento costretto a trovare schemi in un mondo caotico. Meglio vedere schemi dove non ci sono che non vederli dove invece esistono.

\section{L'intelligenza che emerge dal caos}

Ma se vogliamo davvero capire come comportamenti semplici generino \index[main]{complessità emergente}complessità, dobbiamo guardare oltre i singoli animali e osservare cosa accade quando molti animali ``stupidi'' si mettono insieme. Prendiamo gli \index[main]{stormi di uccelli}stormi di uccelli. Avete mai visto migliaia di \index[main]{storni}storni volare insieme al tramonto, disegnando nel cielo quelle forme fluide che sembrano una sola creatura vivente? Si chiama \index[main]{murmuration}murmurazione, e per secoli ha incantato poeti e scienziati. Come fanno migliaia di uccelli a muoversi all'unisono senza scontrarsi, senza un capo, senza un piano?

La risposta è sorprendentemente semplice. Ogni uccello segue soltanto tre regole.\footnote{Reynolds, C. W. (1987). Flocks, herds and schools: A distributed behavioral model. \textit{ACM SIGGRAPH Computer Graphics}, 21(4).} La prima è la separazione: non avvicinarti troppo ai vicini, per evitare le collisioni. La seconda è l'allineamento: vola nella stessa direzione generale di chi ti sta accanto. La terza è la coesione: resta vicino al gruppo, non perderti. Tre regole semplicissime. Nessun uccello comprende la coreografia complessiva, nessuno sa di star creando pattern spettacolari. Eppure da queste tre regole locali emerge un comportamento globale di straordinaria complessità e bellezza. È ciò che i matematici chiamano \index[main]{comportamento emergente}comportamento emergente: il tutto è molto più della somma delle parti.

\begin{center}
	\includegraphics[width=0.9\textwidth]{images/tre_regole.png}
	\captionof{figure}{La dinamica dello stormo di uccelli governata dalle sue tre semplici regole locali.}
	\label{fig:archivista_cervello}
\end{center}

Lo stesso principio governa le formiche. Una formica, da sola, è diciamolo pure piuttosto stupida: il suo cervello ha circa duecentocinquantamila neuroni, contro i nostri ottantasei miliardi. Non sa dove sta andando, non ha un piano, e se la separate dalla colonia probabilmente morirà entro poche ore. Eppure le colonie di formiche costruiscono città sotterranee con sistemi di ventilazione, coltivano funghi, allevano afidi come bestiame, conducono guerre organizzate.\footnote{Hölldobler, B., e Wilson, E. O. (1990). \textit{The Ants}. Harvard University Press, Cambridge.} Come ci riescono? Anche qui, regole semplici che generano complessità. Quando una formica trova del cibo torna al nido lasciando una traccia di \index[main]{feromoni}feromoni; le altre la seguono, e se trovano cibo la rinforzano. Le tracce verso fonti abbondanti vengono rinforzate di più, quelle verso fonti scarse si indeboliscono. Risultato: la colonia trova sempre il percorso più efficiente, senza che nessuna formica abbia mai visto una mappa o calcolato una distanza. È talmente efficiente che gli informatici hanno copiato il sistema creando gli algoritmi delle colonie di formiche, oggi usati per ottimizzare di tutto, dalle rotte di consegna ai circuiti dei microprocessori.\footnote{Dorigo, M., e Stützle, T. (2004). \textit{Ant Colony Optimization}. MIT Press, Cambridge.}

\section{Quando l'intelligenza collettiva diventa stupidità di massa}

Le stesse regole che generano intelligenza collettiva possono produrre \index[main]{stupidità collettiva}stupidità di massa quando il contesto cambia. Pensate a un branco di \index[main]{gnu}gnu che attraversa un fiume. Il primo gnu valuta il punto migliore per guadare, magari sbaglia ma ci prova. Il secondo pensa: ``Lui sa qualcosa che io non so, lo seguo''. Il terzo pensa lo stesso. In pochi minuti migliaia di gnu si gettano nello stesso punto, anche se è palesemente il peggiore, dando origine a quelle scene tragiche di annegamenti di massa che vediamo nei documentari.

È esattamente ciò che accade nelle \index[main]{bolle speculative}bolle speculative. Durante la \index[main]{bolla dei tulipani}bolla dei tulipani del 1637, in Olanda, un singolo bulbo arrivò a costare quanto una casa.\footnote{Mackay, C. (1841). \textit{Extraordinary Popular Delusions and the Madness of Crowds}. Richard Bentley, Londra.} Razionale? No. Ma ogni investitore pensava: ``Se tutti comprano, devono sapere qualcosa che a me sfugge''. Lo stesso \index[cognitive]{bias del gregge}bias del gregge che aiuta gli gnu a trovare l'acqua nella savana li spinge a buttarsi tutti insieme nel posto sbagliato. E noi umani abbiamo perfezionato questa capacità di trasformare l'intelligenza collettiva in stupidità di massa: durante la \index[main]{crisi finanziaria del 2008}crisi finanziaria del 2008, banchieri intelligentissimi, con dottorati nelle migliori università, hanno seguito il gregge comprando titoli che non capivano, semplicemente perché ``lo stavano facendo tutti''. Il nostro cervello sociale, calibrato per gruppi di cinquanta persone nella savana, deve oggi gestire segnali provenienti da milioni di sconosciuti.

\section{Lo specchio peloso: i primati e i nostri bias condivisi}

Se volete vedere i vostri bias cognitivi in azione pura, senza le complicazioni della cultura e del linguaggio, osservate uno scimpanzé o un \index[main]{bonobo}bonobo. Condividiamo con loro il 98\% del DNA, e si vede nei test di laboratorio. Prendiamo l'\index[cognitive]{effetto ancoraggio}effetto ancoraggio. Quando a degli scimpanzé viene mostrato per caso il numero otto prima di scegliere tra diverse quantità di cibo, tendono a preferire quantità vicine all'otto; se vedono il numero tre, scelgono quantità vicine al tre.\footnote{Beran, M. J. (2008). Monkeys track, enumerate, and compare multiple sets of moving items. \textit{Journal of Experimental Psychology: Animal Behavior Processes}, 34(1).} Il numero non ha nulla a che fare con il cibo, eppure ne influenza la scelta. È lo stesso meccanismo per cui, se in un negozio vedete prima una giacca da mille euro, quella da cinquecento vi sembrerà economica. Non è cultura né educazione: è hardware neurologico che condividiamo con i nostri cugini pelosi.

Ma il fenomeno più impressionante è l'avversione alle perdite nei primati. Laurie Santos, a Yale, ha condotto esperimenti brillanti con scimmie cappuccine addestrate a usare gettoni come denaro.\footnote{Chen, M. K., Lakshminarayanan, V., e Santos, L. R. (2006). How basic are behavioral biases? Evidence from capuchin monkey trading behavior. \textit{Journal of Political Economy}, 114(3).} Le scimmie mostrano esattamente gli stessi schemi irrazionali degli umani: preferiscono un guadagno sicuro ma piccolo a uno incerto ma maggiore, eppure quando si tratta di perdite diventano improvvisamente amanti del rischio. In concreto, una scimmia preferisce ricevere con certezza due acini d'uva piuttosto che avere il cinquanta per cento di probabilità di riceverne quattro; ma se deve scegliere tra perdere con certezza due acini o avere il cinquanta per cento di probabilità di perderne quattro, all'improvviso diventa una giocatrice d'azzardo. Stesso cervello, stessa matematica sbagliata, stesso bias che spinge gli umani a tenere troppo a lungo le azioni in perdita sperando in una ripresa.

\section{Il pregiudizio è più antico dell'umanità}

Quella sottile ansia prima di una riunione con il capo, la deferenza automatica verso una figura autorevole, la competizione per lo status sui social: tutto questo non nasce dalla cultura moderna. È l'eco di un imperativo biologico ancora più antico del tribalismo, la gestione delle gerarchie sociali. Gli scimpanzé mostrano quello che possiamo solo chiamare un \index[main]{razzismo}razzismo primitivo: discriminano sistematicamente gli scimpanzé di altri gruppi, anche quando non c'è competizione per le risorse.\footnote{Goodall, J. (1986). \textit{The Chimpanzees of Gombe: Patterns of Behavior}. Harvard University Press, Cambridge.} In esperimenti controllati aiutano più volentieri i membri del proprio gruppo, condividono più cibo con chi assomiglia loro, sono più aggressivi verso chi è diverso. E non è apprendimento: cuccioli cresciuti in isolamento mostrano le stesse preferenze quando vengono esposti a immagini di scimpanzé del proprio gruppo rispetto a estranei. È lo stesso \index[main]{tribalismo}tribalismo che vediamo amplificato negli umani, solo che noi l'abbiamo elevato a forma d'arte: discriminiamo per razza, religione, nazionalità, partito politico, squadra di calcio, marca di telefono. Il meccanismo neurologico è identico, solo le manifestazioni sono diventate barocche.

\section{L'egoismo dell'altruista: un istinto animale}

Ma forse il bias più profondo e controintuitivo che condividiamo con i nostri cugini animali è quello che ci fa credere capaci di \index[main]{altruismo puro}altruismo puro. È un paradosso che ha tormentato i biologi evolutivi per decenni: perché un individuo dovrebbe sacrificare le proprie risorse, o perfino la vita, per un altro? Dal punto di vista della selezione naturale individuale non avrebbe senso, eppure vediamo comportamenti altruistici ovunque nel regno animale. Una delle spiegazioni più influenti suggerisce qualcosa di provocatorio: molto di ciò che chiamiamo altruismo potrebbe avere radici evolutive ``egoistiche''. Non nel senso volgare del termine, ma nel senso che comportamenti apparentemente generosi hanno spesso favorito la sopravvivenza dei nostri geni.

Il primo meccanismo è la \index[main]{selezione parentale}selezione parentale, teorizzata da William Hamilton negli anni Sessanta e resa celebre da Richard Dawkins ne \textit{Il gene egoista}.\footnote{Dawkins, R. (1976). \textit{The Selfish Gene}. Oxford University Press, Oxford.} L'idea è semplice ma rivoluzionaria: quando aiuto i miei parenti, sto in realtà aiutando copie dei miei stessi geni che vivono nei loro corpi. Una delle conferme più convincenti arriva da un luogo inaspettato, le montagne del Tibet, dove vivono piccoli uccelli osservati per anni dai ricercatori che hanno annotato ogni atto di generosità tra parenti. Gli uccelli sembrano seguire una regola matematica precisa, quella che Hamilton aveva teorizzato: il grado di parentela moltiplicato per il beneficio di chi riceve aiuto deve superare il costo sostenuto da chi lo offre.\footnote{Wang, J., e altri (2024). Hamilton's inclusive fitness maintains heritable altruism polymorphism. \textit{Proceedings of the National Academy of Sciences}, 121(8).} Le \index[main]{api operaie}api operaie offrono l'esempio più estremo: quando un'ape punge per difendere l'alveare, muore. Sembra il massimo del sacrificio, ma a causa del loro particolare sistema riproduttivo le operaie condividono con le sorelle il settantacinque per cento dei geni, più di quanto ne condividerebbero con figli propri.\footnote{Foster, K. R., Wenseleers, T., e Ratnieks, F. L. (2006). Kin selection is the key to altruism. \textit{Trends in Ecology and Evolution}, 21(2).} Morire per l'alveare è, dal punto di vista dei geni, un ottimo investimento.

Il secondo meccanismo è l'\index[main]{altruismo reciproco}altruismo reciproco, teorizzato da Robert Trivers nel 1971.\footnote{Trivers, R. L. (1971). The evolution of reciprocal altruism. \textit{The Quarterly Review of Biology}, 46(1).} Qui la logica è diversa: ti aiuto oggi nell'aspettativa, spesso inconscia, che tu aiuterai me domani. Funziona solo quando gli individui si incontrano ripetutamente e ricordano chi ha cooperato e chi ha tradito. I \index[main]{pipistrelli vampiro}pipistrelli vampiro ne sono l'esempio più affascinante: devono nutrirsi di sangue ogni due o tre giorni, altrimenti muoiono, e quando uno torna affamato al dormitorio gli altri rigurgitano sangue per lui. Ma non è carità indiscriminata: tengono nota meticolosa di chi ha condiviso in passato e ricambiano di conseguenza.\footnote{Carter, G. G., e Wilkinson, G. S. (2013). Food sharing in vampire bats: reciprocal help predicts donations more than relatedness or harassment. \textit{Proceedings of the Royal Society B}, 280(1753).} Iniziano persino con investimenti piccoli, pulendosi a vicenda, per saggiare l'affidabilità del partner prima di passare alla costosa condivisione del sangue.\footnote{Carter, G. G., e altri (2020). Development of new food-sharing relationships in vampire bats. \textit{Current Biology}, 30(7).} Anche i suricati, famosi per le loro sentinelle, non sono gli eroi che sembrano: le ricerche mostrano che fanno la guardia da posizioni sicure e sopraelevate, godono dei migliori vantaggi di fuga e sono di solito individui già ben nutriti.\footnote{Clutton-Brock, T. H., e altri (1999). Selfish sentinels in cooperative mammals. \textit{Science}, 284(5420).}

\section{Il cervello altruista: quando aiutare fa sentire bene}

Ma perché ci sentiamo così bene quando aiutiamo gli altri? Le neuroscienze moderne hanno svelato il trucco evolutivo: l'altruismo attiva intensamente i \index[main]{centri di ricompensa}centri di ricompensa del cervello. Studi con risonanza magnetica funzionale mostrano che, quando doniamo a una causa benefica, si accendono le stesse aree, il \index[main]{nucleo accumbens}nucleo accumbens e l'\index[main]{area tegmentale ventrale}area tegmentale ventrale, che si attivano quando mangiamo cioccolato.\footnote{Moll, J., e altri (2006). Human fronto-mesolimbic networks guide decisions about charitable donation. \textit{Proceedings of the National Academy of Sciences}, 103(42).} L'\index[main]{ossitocina}ossitocina, spesso chiamata ormone dell'amore, gioca un ruolo cruciale, ma con un dettaglio rivelatore: aumenta l'altruismo in modo selettivo verso il proprio gruppo. In esperimenti sull'uomo, somministrata per via nasale, fa crescere le donazioni a enti sociali ma riduce quelle a cause ambientali.\footnote{Marsh, N., e altri (2021). Oxytocin-enforced in-group conformity reduces deviant behavior. \textit{Scientific Reports}, 11.} Non è amore universale: è tribalismo biochimico. L'euforia che segue gli atti generosi, quel \index[main]{helper's high}piacere di chi aiuta, non è dunque un sottoprodotto accidentale, ma il modo in cui l'evoluzione ci ricompensa per comportamenti che, statisticamente, hanno favorito i nostri geni. Anche la \index[main]{reputazione}reputazione conta: le persone sono molto più generose quando i loro gesti sono pubblici, perché in una specie sociale come la nostra una buona fama di cooperatore affidabile vale oro.\footnote{Yoeli, E., e altri (2013). Powering up with indirect reciprocity in a large-scale field experiment. \textit{Proceedings of the National Academy of Sciences}, 110(supplemento 2).}

\section{La trappola della teoria della patina: siamo davvero solo scimmie violente?}

Analizzando questa eredità, è facile cadere in una trappola intellettuale tanto seducente quanto pericolosa: l'idea che la civiltà e la moralità siano solo una sottile patina, una vernice fragile stesa su una natura umana di fondo egoista e violenta. È la visione che il primatologo \index[main]{de Waal, Frans}Frans de Waal ha battezzato, per criticarla, \index[main]{teoria della patina}teoria della patina.\footnote{de Waal, F. (2006). \textit{Primates and Philosophers: How Morality Evolved}. Princeton University Press, Princeton.} De Waal, e con lui decenni di ricerca etologica, dimostra che la cooperazione e l'empatia non sono un'invenzione recente, ma affondano nelle radici della nostra linea evolutiva: gli scimpanzé si riconciliano dopo i conflitti, i bonobo sciolgono le tensioni con il sesso anziché con la violenza, le cappuccine protestano contro le ingiustizie nella distribuzione del cibo. L'altruismo non è una vernice, è parte del nostro hardware ancestrale.

Questa, però, non è una disputa nata ieri. È la riedizione moderna, con i dati al posto delle congetture, di una delle più celebri controversie della filosofia politica. Nel 1651 Thomas Hobbes, nel \textit{Leviatano}, descrisse lo stato di natura come una guerra di tutti contro tutti, in cui l'uomo è lupo per l'altro uomo.\footnote{Hobbes, T. (1651). \textit{Leviathan}. Andrew Crooke, Londra.} Per Hobbes, senza un potere assoluto a tenerci a freno ci sbraneremmo a vicenda, e la vita sarebbe solitaria, povera, brutale e breve. Poco più di un secolo dopo, nel 1755, Jean-Jacques Rousseau rovesciò la prospettiva nel suo \textit{Discorso sull'origine della disuguaglianza}: l'uomo allo stato di natura, il famoso buon selvaggio, non era affatto un predatore, ma una creatura sostanzialmente pacifica, e fu semmai la civiltà, con la proprietà privata e le istituzioni, a corromperlo e a renderlo competitivo e crudele.\footnote{Rousseau, J. J. (1755). \textit{Discours sur l'origine et les fondements de l'inégalité parmi les hommes}. Marc-Michel Rey, Amsterdam.} La neuroscienza, oggi, suggerisce una risposta più sfumata di entrambe: Hobbes aveva torto a vedere solo il lupo, perché l'empatia è biologicamente antica quanto l'aggressività; ma anche il buon selvaggio integrale di Rousseau è un'idealizzazione, perché il tribalismo e la violenza verso l'estraneo sono altrettanto radicati. La verità è che portiamo dentro entrambe le possibilità, ed è il contesto a decidere quale emerga.

È qui che la grande intuizione di Rousseau si rivela più feconda di quanto lui stesso potesse documentare. L'idea che siamo ``solo cacciatori violenti'' ignora infatti l'evidenza più schiacciante: \textit{Homo sapiens} è sopravvissuto e ha prosperato non grazie alla forza individuale, ma a una capacità di cooperazione estrema. La crescita collettiva dei figli, la divisione specializzata del lavoro, la condivisione del cibo anche con individui non imparentati sono stati il nostro vero superpotere. I Neanderthal erano fisicamente più forti e avevano cervelli più grandi, eppure li abbiamo surclassati grazie a reti sociali più ampie e a una cooperazione più sofisticata. Su questo punto un altro pensatore, spesso dimenticato, aveva visto lontano: il naturalista e teorico russo Pyotr Kropotkin, che nel 1902 pubblicò \textit{Il mutuo appoggio}.\footnote{Kropotkin, P. (1902). \textit{Mutual Aid: A Factor of Evolution}. William Heinemann, Londra.} Osservando la fauna delle steppe siberiane, Kropotkin contestò la lettura puramente competitiva del darwinismo allora dominante e sostenne che, accanto alla lotta per l'esistenza, esiste un fattore evolutivo altrettanto potente, e spesso più decisivo: l'aiuto reciproco. Le specie che cooperano, argomentava, sopravvivono meglio di quelle che si limitano a competere. Per oltre un secolo questa tesi fu considerata una generosa ingenuità; oggi la biologia evolutiva, con la selezione parentale e l'altruismo reciproco, gli ha dato in larga parte ragione. La cooperazione non è l'eccezione gentile in un mondo di artigli, è uno dei motori centrali dell'evoluzione.

Se però non siamo ``naturalmente'' solo egoisti, perché le nostre società sono così spesso dominate dalla ricerca del potere e dal conflitto? La risposta non sta tanto nella nostra natura quanto nelle strutture che abbiamo costruito. Il potere non è l'unico interesse umano, ma è quello che tende a prevalere nei sistemi sociali su larga scala, non perché siamo intrinsecamente dominatori, ma perché i sistemi stessi, politici ed economici, tendono a selezionare i tratti più spregiudicati. Chi desidera il potere sopra ogni altra cosa ha più probabilità di ottenerlo rispetto a chi valorizza la connessione, la creatività o la cura degli altri. Ridurre ogni interazione umana a una lotta per il potere è intellettualmente comodo ma empiricamente falso: ignora le amicizie che attraversano le classi sociali, la cooperazione che nasce senza coercizione, i sacrifici che nessuna teoria dei giochi riesce a spiegare, e lascia fuori l'arte, l'amore e la curiosità come motivazioni autonome.

Ed è qui che la posizione di Rousseau ritorna in una forma aggiornata e quasi rovesciata. Forse la vera patina è quella al contrario. La barbarie non sarebbe la nostra natura profonda, ma la vernice imposta \textit{sopra} le nostre tendenze cooperative naturali da sistemi di potere che traggono vantaggio dal conflitto. I potenti hanno tutto l'interesse a farci credere che siamo per natura violenti e inaffidabili: così accettiamo la loro violenza e il loro controllo come inevitabili, anziché riconoscerli come costruiti. Riconoscere la nostra profonda eredità cooperativa, allora, non è ingenuità: è quasi un atto di resistenza. Rousseau, che attribuiva alle istituzioni la corruzione di una bontà originaria, aveva intuito proprio questo meccanismo, sia pure con il linguaggio del suo secolo.

C'è infine una conseguenza filosofica che vale la pena cogliere, perché chiude il cerchio. Se la moralità è davvero un prodotto dell'evoluzione, plasmata dalla selezione parentale e dalla reciprocità, ha ancora senso parlare di valori giusti in sé, e non solo di istinti utili alla sopravvivenza dei geni? È la domanda che il filosofo Peter Singer affronta nel suo \textit{L'espansione del cerchio}.\footnote{Singer, P. (1981). \textit{The Expanding Circle: Ethics and Sociobiology}. Farrar, Straus and Giroux, New York.} La sua risposta è sorprendentemente ottimista. È vero, riconosce Singer, che la nostra empatia nasce ristretta: un cerchio piccolo che comprende la famiglia, poi la tribù, poi al massimo la nazione, esattamente come l'ossitocina che premia la generosità solo verso il proprio gruppo. Ma l'uomo possiede uno strumento che gli altri animali non hanno, la ragione, ed è proprio la ragione che ci permette di accorgerci di un'incoerenza: se il dolore di mio fratello conta, perché non dovrebbe contare quello di uno straniero, che lo prova esattamente allo stesso modo? Una volta avviato, questo ragionamento spinge il cerchio dell'empatia ad allargarsi oltre i confini biologici, fino a includere altri popoli, le generazioni future, perfino le altre specie. La storia morale dell'umanità, l'abolizione della schiavitù, l'estensione dei diritti, è per Singer proprio questo lento ampliamento del cerchio. La nostra natura ci consegna un punto di partenza ristretto, ma non ci condanna a restarvi: la stessa intelligenza che ci ha resi tribali può renderci, faticosamente, universali.

\begin{figure}[htbp]
	\centering
	\includegraphics[width=0.8\textwidth]{images/L_Espansione_del_Cerchio_Morale (1).png}
	\captionof{figure}{L'espansione del cerchio morale secondo Peter Singer: l'empatia nasce ristretta alla famiglia e alla tribù, dove agiscono la selezione parentale e la reciprocità, ma la ragione ci permette di allargarla progressivamente fino a includere l'intera umanità, le altre specie e le generazioni future.}
	\label{fig:cerchio_morale}
\end{figure}

\section{Le gerarchie dei primati e l'ansia dell'ufficio}

C'è un altro aspetto della nostra eredità primordiale che merita attenzione: la gestione delle \index[main]{gerarchie sociali}gerarchie sociali. Negli scimpanzé, come in quasi tutte le società di primati, esistono \index[main]{gerarchie di dominanza}gerarchie di dominanza molto rigide, e ogni individuo sa esattamente qual è il suo posto. Le neuroscienze ci dicono che non si tratta solo di comportamento appreso, ma di chimica cerebrale: i livelli di \index[main]{serotonina}serotonina, il neurotrasmettitore associato al benessere e alla fiducia in sé, sono molto più alti negli individui alfa e crollano se questi perdono il loro status.\footnote{Raleigh, M. J., e McGuire, M. T. (1994). Serotonin, aggression, and violence in vervet monkeys. In \textit{The Neurotransmitter Revolution}. MIT Press, Cambridge.} Gli individui di rango inferiore si sottomettono spontaneamente, mostrano segni fisiologici di stress, con livelli più alti di \index[main]{cortisolo}cortisolo, ed evitano il contatto visivo diretto con i dominanti.\footnote{Sapolsky, R. M. (2017). \textit{Behave: The Biology of Humans at Our Best and Worst}. Penguin Press, New York.}

Vi suona familiare? Dovrebbe. È esattamente ciò che accade in ogni ufficio, ogni scuola, ogni organizzazione umana. Frans de Waal ha descritto le sale riunioni delle aziende come un'arena per le stesse strategie di coalizione e dominanza osservate per decenni nelle sue colonie di scimpanzé.\footnote{de Waal, F. (1982). \textit{Chimpanzee Politics: Power and Sex among Apes}. Johns Hopkins University Press, Baltimora.} Abbiamo gli stessi schemi di sottomissione e dominio, gli stessi segnali non verbali, la stessa \index[main]{ansia gerarchica}ansia gerarchica. Il celebre \index[main]{Whitehall Study}studio Whitehall sugli impiegati statali britannici ha dimostrato una correlazione sconcertante: più basso è il grado nella gerarchia lavorativa, più alto è il rischio di malattie da stress, come gli infarti.\footnote{Marmot, M. G., e altri (1991). Health inequalities among British civil servants: the Whitehall II study. \textit{The Lancet}, 337(8754).} Il nostro corpo reagisce a una email critica del capo con lo stesso allarme fisiologico con cui un primate di basso rango reagisce allo sguardo minaccioso dell'alfa. La differenza è che noi abbiamo inventato gerarchie infinitamente più astratte: mentre gli scimpanzé competono per cibo e partner, noi competiamo per stipendi, titoli, seguito sui social, la marca dell'auto o il quartiere in cui viviamo. Abbiamo semplicemente moltiplicato i modi simbolici di giocare allo stesso antichissimo gioco dello \index[main]{status sociale}status sociale, ma il software neurobiologico che lo gestisce è rimasto identico.

\section{La saggezza e la follia delle folle}

C'è un esperimento celebre, condotto per la prima volta da Francis Galton nel 1906.\footnote{Galton, F. (1907). Vox populi. \textit{Nature}, 75(1949).} A una fiera di paese, ottocento persone dovevano indovinare il peso di un bue. Prese singolarmente, le stime erano pessime. Ma la media di tutte le stime risultò di circa cinquecentoquarantatré chili, contro un peso reale di poco superiore: un errore inferiore allo zero virgola uno per cento. È la \index[main]{saggezza delle folle}saggezza delle folle: in certe condizioni i gruppi sono più intelligenti dei singoli individui più intelligenti del gruppo. Ma, ed è cruciale, solo quando ciascuno formula la propria valutazione in modo indipendente. Appena le persone possono influenzarsi a vicenda, la saggezza collettiva collassa. È esattamente ciò che accade nei \index[main]{social media}social media: invece di ottocento persone che stimano in autonomia il peso di un bue, abbiamo milioni di individui che si influenzano l'un l'altro su vaccini, politica ed economia. Il risultato non è saggezza collettiva ma \index[main]{polarizzazione}polarizzazione, con gruppi che diventano più estremi delle opinioni di partenza.

\section{Il paradosso della nostra epoca}

Eccoci dunque al \index[main]{paradosso}paradosso del nostro tempo. Siamo \index[main]{animali sociali}animali sociali con cervelli progettati per gruppi di cinquanta individui, e viviamo in un mondo di otto miliardi di persone interconnesse. Siamo topi superstiziosi con armi nucleari, formiche che seguono tracce di feromoni digitali chiamati algoritmi, scimpanzé tribali con un account sui social. Il nostro problema non è avere questi bias, che sono anzi la ragione per cui siamo sopravvissuti come specie: il problema è che li stiamo applicando in un contesto radicalmente diverso da quello per cui si sono evoluti. È come usare una mappa del Cinquecento per navigare con il GPS: alcuni punti di riferimento sono ancora validi, ma le strade sono cambiate del tutto.

C'è però anche speranza in questa storia. Se le formiche possono costruire città senza architetti e gli storni danzare nel cielo senza coreografi, forse anche noi possiamo imparare a trasformare i nostri bias da debolezze in forze, non eliminandoli, perché sono cablati nel cervello, ma comprendendoli e costruendo sistemi che ne tengano conto. \index[main]{Wikipedia}Wikipedia funziona perché sfrutta insieme il nostro desiderio di contribuire e la nostra pignoleria nel correggere gli errori altrui. La scienza stessa funziona non nonostante i bias degli scienziati, ma perché ha sviluppato meccanismi, la revisione tra pari, la replicazione, la falsificazione, che li correggono nel tempo. È, in fondo, l'espansione del cerchio di cui parlava Singer applicata alla conoscenza: strumenti collettivi che ci permettono di andare oltre i limiti del singolo cervello ancestrale.

\section{La lezione degli animali}

Alla fine, cosa ci insegnano i topi superstiziosi, le formiche organizzate e i nostri cugini scimpanzé? Che l'\index[main]{irrazionalità}irrazionalità non è un difetto del sistema: è il sistema. Che l'\index[main]{intelligenza}intelligenza non è calcolo perfetto, ma adattamento imperfetto. Che siamo tutti, dal topo al primate, dalla formica all'umano, creature che navigano l'incertezza con strumenti cognitivi imperfetti eppure sorprendentemente efficaci. La continuità evolutiva dei bias attraverso le specie ci mostra che non siamo angeli della ragione caduti in corpi animali: siamo animali che hanno evoluto la capacità di riflettere sulla propria animalità. Un topo non sa di essere superstizioso, una formica non sa di seguire un algoritmo, uno scimpanzé non sa di essere tribale. Noi lo sappiamo. Non ci rende immuni, ma ci rende consapevoli. E la consapevolezza, pur non essendo una cura, è il primo passo verso la saggezza. Forse la vera intelligenza non sta nel superare la nostra natura animale, ma nell'imparare a danzare con essa, come gli storni che trasformano tre semplici regole in una coreografia di pura bellezza.

\section{Il motore sotto il cofano}

Questa continuità evolutiva dei bias, dal topo all'umano, ci lascia con una domanda di fondo: qual è il linguaggio comune che guida tutti questi comportamenti istintivi? Qual è il motore che spinge il topo a fuggire, la formica a seguire una traccia, lo scimpanzé a diffidare dell'estraneo e l'umano a sacrificarsi, in apparenza, per gli altri? La risposta non sta nella logica o nel calcolo, ma in qualcosa di molto più antico e potente: le \index[main]{emozioni}emozioni. Sono il vero sistema operativo ancestrale, il software primordiale che gira incessantemente sotto la superficie della coscienza. Un topo non pensa di avere paura, semplicemente ha paura; uno scimpanzé non decide di adirarsi, semplicemente si arrabbia; e noi non calcoliamo quando essere altruisti, sentiamo l'impulso di aiutare, guidati da meccanismi che ci ricompensano per comportamenti utili ai nostri geni.

Abbiamo ereditato questa stessa architettura. I nostri bias non sono semplici errori di calcolo, sono le manifestazioni di un motore potente ma obsoleto, progettato per la savana e oggi costretto a navigare nel traffico digitale del ventunesimo secolo. Per capire davvero la nostra vulnerabilità nell'era moderna non basta osservare il comportamento in superficie: dobbiamo aprire il cofano e guardare quel motore da vicino. Nel prossimo capitolo faremo esattamente questo. Esploreremo come i nostri circuiti primitivi, la ricerca di novità, l'attenzione al pericolo, il desiderio di gratificazione, siano stati traditi dall'evoluzione e vengano oggi sfruttati sistematicamente dall'ecosistema digitale. Scopriremo che la tecnologia moderna non è un semplice strumento, ma un ambiente progettato per dialogare direttamente con il nostro motore ancestrale, spesso scavalcando del tutto il pensiero razionale. Preparatevi a scoprire l'anatomia della dipendenza digitale e, soprattutto, a costruire un kit di sopravvivenza per riprendere il controllo.